% ============================================================
%  THE ORTYX PROTOCOL
%  Part V — Methodological Synthesis
%  Chapter 21: The Protocol Named
% ============================================================

\chapter{The Protocol Named}
\label{ch:protocol-named}

\chapepigraph{A methodology that has been developing throughout
a book without being named is not coy. It is being honest
about the order of operations. You cannot name a procedure
before you know what it requires. You only know what it
requires after you have tried to apply it.}{Flyxion}

\noindent
The \Ortyx{} Protocol has been operating since Chapter~6
without having been named as such. It was named in Chapter~11
when the pressure generated by Parts~I and~II made its
formal statement necessary. Now, having been applied to
the Phoenix corpus, having reconstructed the survivable
components in RSVP-compatible terms, and having introduced
the Nexus as a parallel case study, the protocol can be
stated in its fullest form.

Part~V does three things. Chapter~21 states the protocol
explicitly as a general-purpose epistemic methodology,
not a Phoenix-specific tool. Chapter~22 applies the
protocol recursively to itself, asking whether Ortyx
passes its own audit. Chapter~23 addresses the
civilizational context: why a protocol of this kind
is needed now, when it was not needed in earlier
intellectual periods, and what its limits are. Chapter~24
operationalizes the protocol in computational systems,
connecting it to the constraint-safe engineering
principles that the Nexus chapter introduced.

\section{The Protocol as General Methodology}
\label{sec:protocol-general}

The \Ortyx{} Protocol is a general-purpose procedure
for auditing high-coherence speculative frameworks ---
frameworks that exhibit local explanatory power,
architectural elegance, and vocabulary continuity in
ways that make them difficult to evaluate from inside
their own conceptual structure. It applies whenever
the question ``Is this framework correct?'' is preceded
by the more fundamental question ``Is this framework
structured in a way that allows the question of its
correctness to be answered?''

The protocol operates in five phases, corresponding
to the five parts of this monograph:

\textbf{Phase 1: Reconstruction.} The framework is
encountered on its own terms and reconstructed at its
strongest. The auditor suspends judgment and follows
the framework's arguments as sympathetically as possible.
The goal is to understand what the framework's
practitioners find compelling about it, and to produce
a reconstruction that its practitioners would recognize
as fair.

\textbf{Phase 2: Tension Accumulation.} The auditor
begins noticing asymmetries without yet naming them.
Vocabulary migration is observed. Accommodation
strategies are noted. The boundary between genuine
generalization and atmospheric expansion begins to
become visible. No formal critique is mounted; the
observations are recorded informally.

\textbf{Phase 3: Formal Decomposition.} The four
audit operators are applied. The four failure geometries
are identified and their severity assessed. The four-level
taxonomy is applied to classify claims by type. RSVP
or an equivalent reinterpretive framework is audited.
The decomposition produces $(\ThetaS, \ThetaU)$.

\textbf{Phase 4: Reconstruction.} The survivable
components $\ThetaS$ are translated via the reinterpretation
functor $\RSVPfunctor$ into a more stable expressive
context $\ThetaStar$. New case studies consistent with
$\ThetaStar$ are introduced if available. The reconstructed
paradigm is stated as a research programme with explicit
hard core and protective belt.

\textbf{Phase 5: Recursive Self-Application.} The
protocol is applied to itself. The auditor asks which
of the four failure geometries the protocol itself
exhibits, what its own projection dependencies are,
what would constitute evidence against its core claims,
and how it would handle a framework that the protocol
was unable to classify. The self-application produces
$\OrtOp_{n+1} = \AuditOp(\OrtOp_n)$.

\begin{protocol}
\label{prot:ortyx}
The \Ortyx{} Protocol is the five-phase procedure
described above. Its application to a framework
$\SysTriple = (\mathcal{A}, \mathcal{M}, \mathcal{E})$
produces:
\[
  \text{Phase 3:}\quad
  \SysTriple \xrightarrow{\AuditOp} (\ThetaS, \ThetaU),
\]
\[
  \text{Phase 4:}\quad
  \ThetaS \xrightarrow{\RSVPfunctor} \ThetaStar,
\]
\[
  \text{Phase 5:}\quad
  \OrtOp \xrightarrow{\AuditOp} \OrtOp'.
\]
The protocol is complete when Phase~5 returns a stable
version of the protocol: when $\OrtOp'$ does not
substantially revise the operators, weights, or
procedures of $\OrtOp$. The protocol has not yet
achieved stability in this sense; the present chapter
is part of the iteration.
\end{protocol}

\section{What the Protocol Requires of Its User}
\label{sec:protocol-requirements}

The protocol makes specific demands on the person
applying it that are worth stating explicitly, because
they are easily violated under the pressures of
intellectual engagement.

\textit{Phase 1 requires genuine sympathy.} The
reconstruction must not be a caricature. If the auditor
finds the framework boring, implausible, or obviously
wrong, Phase~1 will not be done well, and the later
phases will be operating on a distorted target. The
discipline required is to find the framework genuinely
interesting for long enough to reconstruct it at its
strongest.

\textit{Phase 2 requires patience.} The temptation
is to move to formal critique as soon as an asymmetry
is noticed. Resisting this temptation is essential:
the informal accumulation of tension gives the formal
critique its legitimacy. A critique that arrives before
the reader has inhabited the framework feels like
external imposition; a critique that arrives after
will feel like internal necessity.

\textit{Phase 3 requires precision.} The failure
geometries are not rhetorical tools; they are formal
categories that require specific evidence for each
application. Asserting $\FailGeo{3}$ (Semantic
Universality Collapse) requires demonstrating that
the relevant concept has been broadened to the point
of accommodating both the presence and absence of
relevant phenomena. It is not sufficient to say that
the framework is vague.

\textit{Phase 4 requires honesty about the reinterpretive
framework.} The functor $\RSVPfunctor$ is not a
compliment to RSVP. It is a tool. The auditor who
finds that the reinterpretive framework has worse audit
scores than the original framework must acknowledge
this and either find a better reinterpretive framework
or revise the conclusion.

\textit{Phase 5 requires courage.} The self-application
is not performative humility. It must find real
vulnerabilities in the protocol, not theatrical ones.
An auditor who applies the protocol to itself and finds
nothing must either have a very good explanation for
why the protocol is uniquely exempt from its own
criteria, or must try harder.

\section{The Scope of the Protocol}
\label{sec:protocol-scope}

The \Ortyx{} Protocol is applicable to any theoretical
framework that: (1)~makes claims at multiple inferential
levels without clearly distinguishing them, (2)~exhibits
vocabulary that migrates across domains carrying authority
from its source domain, (3)~has a core concept that
could in principle undergo Semantic Universality Collapse,
and (4)~makes claims that depend on projection operators
whose constraint-faithfulness is not established.

These conditions are very common. They apply to most
speculative frameworks in cognitive science, consciousness
studies, AI research, and philosophy of mind, as well
as to many frameworks in the social sciences and
humanities where quantitative vocabulary is borrowed
from physics or computer science.

They do not apply to frameworks that: (1)~operate
entirely within a single, well-delimited domain with
established operationalization conventions, (2)~make
no cross-domain claims, and (3)~have vocabulary that
does not migrate. Such frameworks may have other problems
but they are not the target of the \Ortyx{} Protocol.

\begin{auditprop}
\label{ap:scope-limitation}
The \Ortyx{} Protocol has a defined scope. Applying
it to frameworks outside that scope --- frameworks
whose problems are different from the failure geometries
the protocol identifies --- would be a misuse of the
procedure. The protocol is not a general suspicion
machine; it is a tool for a specific class of epistemic
failure. Part of responsible use of the protocol is
recognizing when a framework's problems are not the
problems the protocol addresses.
\end{auditprop}

\section{The Protocol and Generative Conditions}
\label{sec:generative-conditions}

The Preface noted that the conditions that produce
locally compelling but inferentially unstable frameworks
--- symbolic abundance, rapid production cycles,
optimization pressure toward legibility, and generative
amplification of coherent-sounding prose --- now operate
at institutional scale. The \Ortyx{} Protocol was
developed in response to a specific corpus, but the
problem it addresses is structural, not marginal.

The generative conditions that the Preface identified
are worth examining in the light of the completed
protocol. Large language models, operating at scale,
are extraordinarily capable of producing text that
exhibits: vocabulary coherence (the same terms appear
consistently), organizational formalism (complex claims
are expressed in formally organized structures), and
broad explanatory scope (the vocabulary can be applied
to many domains). These are precisely the features
that the \Ortyx{} Protocol identifies as warning signs
when they occur without the corresponding properties
of operational grounding, constraint propagation,
and falsification structure.

This is not an argument that LLM-generated text is
uniformly poor. It is an argument that the tools developed
in this monograph for evaluating the Phoenix corpus
are also useful for evaluating LLM-assisted intellectual
production more broadly. The four failure geometries
are not specific to the Phoenix corpus; they are failure
modes that occur whenever the infrastructure for
producing fluent, organized, vocabulary-coherent text
outpaces the infrastructure for ensuring that the text
provides genuine inferential constraint.

\medskip

Chapter~22 applies the protocol to itself. The analysis
will not be comfortable, because the same tools that
identified the Phoenix corpus's vulnerabilities apply
to the protocol that developed them. The test of
intellectual seriousness is whether the analysis is
conducted honestly rather than performed theatrically.
